% !TEX root = ../../tree_notes.tex
% =====================================================================
%  Chapter 7 opener -- "Chapter at a Glance", reading paths + ladder
%  staircase, the running BST B seen four ways (sorted line, search path,
%  value intervals, container view), roadmap, recognition guide,
%  constraint table, notation.  NO numbered boxes here.
% =====================================================================

\begin{chaptertheorybox}{Chapter at a Glance}
Chapters~2 and~3 walked trees that made \emph{no promises about their contents}. This chapter
installs a \term{rule} inside the shape: at every node, everything in the left subtree is
\term{smaller} and everything in the right subtree is \term{larger}. That single invariant --
$\text{left} < \text{node} < \text{right}$, \emph{recursively} -- converts a binary tree into a
\term{Binary Search Tree} (BST), and a BST is \emph{binary search wearing a tree costume}: the
sorted order you already exploit in an array becomes available \emph{structurally}, as a
root-to-leaf descent. Three consequences run this whole chapter:
\begin{itemize}[nosep]
  \item \textbf{The inorder walk \emph{is} the sorted list.} One traversal property
        (\thmref{c03-one-walk} in \chref{c03}, never re-derived) plus the BST rule, and every
        ``ordered question'' -- kth smallest, closest value, minimum gap -- becomes a walk of a
        sorted line that happens to be stored in a tree. \hfill \lv{2}--\lv{3}
  \item \textbf{Search, insert and delete are one descent.} Compare, go left or right, repeat:
        $\bigO{h}$ steps, exactly the halving intuition of binary search -- and every modification
        must \emph{preserve} the invariant, which is what makes deletion (three cases, successor
        splice) the chapter's classic interview mechanic. \hfill \lv{2}--\lv{4}
  \item \textbf{The height is the whole cost -- and it can rot.} Insert $1,2,\dots,n$ in order and
        the BST degenerates into a chain: every descent becomes $\bigO{n}$ and the data structure
        is a linked list with extra pointers. Every serious implementation fixes this by
        \emph{balancing} -- the promise this chapter makes and Chapter~8 (\chref{c06}) keeps.
        \hfill \lv{3}--\lv{5}
\end{itemize}

\textbf{You will learn to}
\begin{itemize}[nosep]
  \item state and check the BST property with \emph{bounds} (the local-comparison trap that fails
        most first attempts at validation); \hfill\lv{2}--\lv{3}
  \item search, insert, and delete in one root-to-leaf descent, handling the two-child deletion
        via successor splicing; \hfill\lv{2}--\lv{4}
  \item answer ordered queries -- kth smallest, successor, closest value, minimum gap -- by riding
        the inorder line with early exits; \hfill\lv{2}--\lv{3}
  \item rebuild a BST from a single preorder scan using bound descent (no second array needed);
        \hfill\lv{3}
  \item use \code{std::set}/\code{std::map} as a ready-made balanced BST and recognise the contest
        problems (intervals, nearest keys, calendars) that are one \code{lower\_bound} away.
        \hfill\lv{3}--\lv{5}
\end{itemize}

\textbf{Prerequisites.} Binary-tree recursion from Chapter~2 (\chref{c02}) and the three DFS
orders from Chapter~3 (\chref{c03}) -- especially \emph{inorder}, which this chapter uses as the
sorted line without re-deriving it. The rebuilding-from-two-orders machinery is Chapter~3
(\probref{c03-build-in-post}); this chapter is the \emph{one-order} and \emph{ordered-content} version.
\end{chaptertheorybox}

% ---------------- Station 1b: reading paths + ladder position ----------------
\begin{chaptertheorybox}{How to Read This Chapter (two paths up the ladder)}
Level badges $\lv1$--$\lv6$ as before: $\lv2$ is ``descend and compare'', $\lv3$ is the solid
interview medium, $\lv5$ is an LC Hard where the naive BST's height is the enemy.

\textbf{Beginner path ($\lv2$--$\lv3$).} Read \secref{c05-property} fully -- it fixes the BST rule,
the bounds-based validation, and the one-descent search. Then \secref{c05-insert-delete} through the
three deletion cases. Take every $\lv2$--$\lv3$ solved problem. You may \emph{skip on a first
pass}: the preorder-reconstruction theorem of \secref{c05-queries}, the container engineering of
\secref{c05-contests}, every $\lv4{+}$ variation, and the Summit. At the end of this path you can
search, insert, delete, and answer any ordered query in a BST in $\bigO{h}$.

\textbf{Full path.} Everything, in order -- it ends with the scaled Summit (LC Hard 315
\emph{Count of Smaller Numbers After Self}: insert right-to-left into a size-augmented BST, read
each rank off the left-subtree counts, and confront the $\bigO{n^2}$ chain degeneration head-on)
and the reduction gallery showing that ``kth smallest'', ``minimum gap'', ``inorder successor'',
``range count'' and ``is this interval covered'' are all \emph{the same sorted line} queried four
different ways. The balancing machinery that repairs the rotting height -- AVL, red-black, treap
-- is Chapter~8 (\chref{c06}); the array-indexed cousins (Fenwick, segment tree) that answer the
same rank queries faster are Chapters~18--19 (\chref{c14}, \chref{c13}).
\end{chaptertheorybox}

\begin{figure}[H]
\centering
\begin{tikzpicture}[diagram, x=9mm, y=7mm]
  \foreach \st/\lab in {1/lv1,2/lv2,3/lv3,4/lv4,5/lv5,6/lv6} {
    \fill[\lab!22, draw=\lab, line width=0.8pt] (\st-1,0) rectangle (\st,\st*0.52);
    \node[font=\scriptsize\bfseries, text=\lab] at (\st-0.5,\st*0.52-0.42) {L\st};}
  \node[annot, anchor=south west, font=\scriptsize]
    at (0.05,0.62) {7.1 property, bounds, search};
  \node[annot, anchor=south west, font=\scriptsize]
    at (1.05,1.14) {7.1/7.3 ordered queries};
  \node[annot, anchor=south west, font=\scriptsize]
    at (2.05,1.66) {7.2 insert; three-case delete};
  \node[annot, anchor=south west, font=\scriptsize]
    at (3.05,2.18) {7.3 preorder rebuild; 7.4 ordered containers};
  \node[annot, anchor=south west, font=\scriptsize, text=lv5]
    at (4.05,2.70) {7.4 interval engineering; degeneration analysis};
  \node[annot, anchor=south west, font=\scriptsize, text=lv5]
    at (5.05,3.22) {scaled summit (315)};
\end{tikzpicture}
\caption*{Chapter~7 on the six-step ladder: comparing and descending carry you from the bare
invariant to ordered queries ($\lv3$); deletion, reconstruction and container work sit at
$\lv3$--$\lv4$; the Summit (315) makes the \emph{height itself} the problem, at $\lv5$. Beginner
path $=$ steps 1--3.}
\end{figure}

% ---------------- The running example ----------------
\subsection*{The Running Example: One BST, Four Ways of Seeing It}

This chapter has \emph{one} running example: the BST $B$ below, drawn once and then read four
different ways. Every topic revisits it -- searching and validating (7.1), deleting and inserting
(7.2), ordered queries and a preorder rebuild (7.3), and finally the same contents living inside a
\code{std::set} (7.4). Keys are distinct integers; the tree holds $n=9$ of them:

\begin{figure}[H]
\centering
\begin{tikzpicture}[diagram]
  \node[troot] (n8)  at (0,0)    {$8$};
  \node[tnode] (n3)  at (-3,-1.4) {$3$};
  \node[tnode] (n10) at (3,-1.4)  {$10$};
  \node[tleaf] (n1)  at (-4.2,-2.8) {$1$};
  \node[tnode] (n6)  at (-1.8,-2.8) {$6$};
  \node[tleaf] (n14) at (3,-2.8)  {$14$};
  \node[tleaf] (n4)  at (-2.7,-4.2) {$4$};
  \node[tleaf] (n7)  at (-0.9,-4.2) {$7$};
  \node[tleaf] (n13) at (1.8,-4.2)  {$13$};
  \draw[tedgehl] (n8) -- (n3);  \draw[tedgehl] (n8) -- (n10);
  \draw[tedgehl] (n3) -- (n1);  \draw[tedgehl] (n3) -- (n6);
  \draw[tedgehl] (n6) -- (n4);  \draw[tedgehl] (n6) -- (n7);
  \draw[tedgehl] (n10) -- (n14); \draw[tedgehl] (n14) -- (n13);
  \node[annot, align=left, font=\scriptsize] at (-6.8,-1.4)
    {$n=9$, height $3$\\ inorder $=$ sorted\\ $\mathit{lo}\le\text{key}\le\mathit{hi}$\\ at every node};
  \node[annot, align=left, font=\scriptsize] at (6.2,-2.0)
    {not balanced:\\ deleting $8$ leaves\\ height $4$ (the rot\\ of 7.4's warning)};
\end{tikzpicture}
\caption{The running BST $B$. Root $8$; every left subtree holds smaller keys, every right subtree
larger keys. Nine keys, eight edges, height $3$ -- and one deletion away from being a chain.
That fragility is the thread to Chapter~8.}\label{fig:c05-running}
\end{figure}

\medskip
\noindent Keep these four ``views'' of $B$ in mind -- each is proved or used, never merely
asserted, in the section named below:
\begin{center}
\renewcommand{\arraystretch}{1.25}
\setlength{\tabcolsep}{5pt}
\begin{tabular}{@{}p{0.24\linewidth} p{0.70\linewidth}@{}}
\toprule
\textbf{View} & \textbf{Value for $B$} \\
\midrule
inorder (the sorted line)   & $1,3,4,6,7,8,10,13,14$ -- sorted \emph{because} $B$ is a BST (\thmref{c05-inorder}) \\
preorder (rebuild input)    & $8,3,1,6,4,7,10,14,13$ -- one scan regenerates $B$ (\secref{c05-queries}) \\
value interval of a subtree & node $3$ holds $[1,7]$; node $6$ holds $[4,7]$; node $14$ holds $[13,14]$ (\thmref{c05-bounds}) \\
\code{std::set} contents    & the same nine keys; \code{lower\_bound(6)} $\to$ iterator at $6$, \code{upper\_bound(8)} $\to$ $10$ (\secref{c05-contests}) \\
\bottomrule
\end{tabular}
\end{center}

% ---------------- Roadmap diagram ----------------
\begin{figure}[H]
\centering
\begin{tikzpicture}[diagram, node distance=5mm,
    top/.style={boxnode, fill=primaryblue!10, draw=primaryblue, text width=25mm,
                minimum height=11mm, font=\small},
    res/.style={annot, text width=25mm, anchor=north}]
  \node[top] (t1) {\textbf{7.1} Property \& Search};
  \node[top, right=of t1] (t2) {\textbf{7.2} Insert \& Delete};
  \node[top, right=of t2] (t3) {\textbf{7.3} Queries \& Orders};
  \node[top, right=of t3] (t4) {\textbf{7.4} In Contests};
  \foreach \a/\b in {t1/t2, t2/t3, t3/t4} \draw[flow] (\a) -- (\b);
  \node[res] at ([yshift=-2mm]t1.south) {inorder $=$ sorted; bounds validation; one-descent search};
  \node[res] at ([yshift=-2mm]t2.south) {insert as failed search; three delete cases; successor splice};
  \node[res] at ([yshift=-2mm]t3.south) {kth smallest; successor; min gap; rebuild from preorder};
  \node[res] at ([yshift=-2mm]t4.south) {std::set/map API; interval unions; the height rot};
\end{tikzpicture}
\caption*{The four topics, left to right, are the \emph{same invariant} exploited at increasing
distance from the metal: 7.1 \emph{reads} the rule; 7.2 \emph{edits} under the rule; 7.3
\emph{queries} the sorted line the rule guarantees; 7.4 outsources the whole structure to
\code{std::set} and learns when that is (and is not) the right call.}
\end{figure}

% ---------------- Topic table ----------------
\begin{center}
\renewcommand{\arraystretch}{1.25}
\setlength{\tabcolsep}{4pt}
\begin{tabular}{@{}p{0.24\linewidth} p{0.36\linewidth} c p{0.23\linewidth}@{}}
\toprule
\textbf{Topic} & \textbf{Key results} & \textbf{Levels} & \textbf{Solved problems} \\
\midrule
\secref{c05-property}\ \ The BST Property \& Searching
  & inorder $=$ sorted both ways; bounds validation; $\bigO{h}$ descent & $\lv2$--\lv{3}
  & \probref{c05-search} \\
\addlinespace
\secref{c05-insert-delete}\ \ Insertion and Deletion
  & insert as failed search; three delete cases; successor splice & $\lv2$--\lv{4}
  & \probref{c05-delete} \\
\addlinespace
\secref{c05-queries}\ \ Queries \& Order Statistics
  & kth/smallest-line rides inorder; min gap; rebuild from preorder alone & $\lv2$--\lv{3}
  & \probref{c05-construct} \\
\addlinespace
\secref{c05-contests}\ \ BSTs in Contests (\code{std::set/map})
  & ordered-container API; interval management; degeneration warning & $\lv3$--\lv{5}
  & \probref{c05-range} \\
\bottomrule
\end{tabular}
\end{center}

% ---------------- Problem recognition guide ----------------
\subsection*{Problem Recognition Guide --- When You See \emph{this}, Think \emph{that}}

BST problems announce themselves by \emph{promising an order}: keys, values, or intervals with a
guaranteed comparison structure. Match the phrasing; every row points at the section that builds
the tool.

\begin{center}
\renewcommand{\arraystretch}{1.3}
\setlength{\tabcolsep}{4pt}
\begin{tabular}{@{}p{0.42\linewidth} p{0.35\linewidth} l@{}}
\toprule
\textbf{When the problem says\ldots} & \textbf{\textcolor{deepblue}{\textbf{Think\ldots}}} & \textbf{Section} \\
\midrule
``each node has $\mathit{left} < \mathit{root} < \mathit{right}$'', ``a BST is given''
  & the whole tree is binary search -- descend, never scan & \secref{c05-property} \\
\addlinespace
``is this a valid BST?'', ``check the ordering''
  & bounds descent, not local child comparisons (\thmref{c05-bounds}) & \secref{c05-property} \\
\addlinespace
``insert a key'', ``delete a node'', ``prune the tree''
  & one descent; for two-child deletion, splice the successor & \secref{c05-insert-delete} \\
\addlinespace
``kth smallest'', ``minimum difference'', ``successor of $p$'', ``closest to target''
  & ride the inorder sorted line with an early exit & \secref{c05-queries} \\
\addlinespace
``given only the preorder, rebuild the BST''
  & bound descent: next node attaches under the last open bound & \probref{c05-construct} \\
\addlinespace
``add/query/remove intervals'', ``nearest free seat'', ``no double booking''
  & \code{std::map/set} \code{lower\_bound} -- an off-the-shelf BST & \secref{c05-contests} \\
\addlinespace
``count smaller elements after self'' (and $n$ large)
  & order statistics per step; watch the naive BST degenerate (Summit) & \probref{c05-summit} \\
\bottomrule
\end{tabular}
\end{center}

\begin{center}
\renewcommand{\arraystretch}{1.25}
\setlength{\tabcolsep}{4pt}
\begin{tabular}{@{}p{0.27\linewidth} p{0.45\linewidth} p{0.20\linewidth}@{}}
\toprule
\textbf{Constraint you see} & \textbf{What it forces} & \textbf{Chapter-7 tool} \\
\midrule
``it is a BST'' in the statement & never traverse all $n$ -- descend $\bigO{h}$ & \secref{c05-property} \\
validation, or any ``closest to $x$'' & carry $(\mathit{lo},\mathit{hi})$ down the descent & \thmref{c05-bounds} \\
delete with two children & successor/predecessor splice, one extra descent & \secref{c05-insert-delete} \\
insertion order chosen by the input (sorted!) & the naive BST becomes a chain -- balance or re-plan & \secref{c05-contests} (\chref{c06}) \\
dense integer keys, many rank queries & Fenwick/segment beats a pointer BST & \chref{c14}, \chref{c13} \\
\bottomrule
\end{tabular}
\end{center}

\begin{notebox}[Not this chapter?]
If the tree's contents carry \emph{no order} and you just need to visit it, that is Chapters~2--3
(\chref{c02}, \chref{c03}). If you must rebuild a binary tree from \emph{two} traversal orders,
or serialise/parse it, that is Chapter~3 (\probref{c03-build-in-post}) -- this chapter needs only
\emph{one} order, because the BST rule does the other half of the work. If you want the
\emph{guarantee} that the height stays $\bigO{\log n}$ -- AVL rotations, red-black colours, treap
priority -- that is Chapter~8 (\chref{c06}); we \emph{assume} containers are balanced and show what
breaks when you cannot. If keys are dense integers and you need ranks \emph{by array index}, the
Fenwick tree of Chapter~18 (\chref{c14}) is the faster engine. This chapter owns \emph{the ordered
tree itself: the rule, the descent, and every question the sorted line answers}.
\end{notebox}

% ---------------- Notation ----------------
\subsection*{Notation Used in This Chapter}

\begin{center}
\renewcommand{\arraystretch}{1.2}
\setlength{\tabcolsep}{4pt}
\begin{tabular}{@{}p{0.20\linewidth} p{0.27\linewidth}@{\quad}p{0.17\linewidth} p{0.27\linewidth}@{}}
\toprule
\textbf{Symbol} & \textbf{Meaning} & \textbf{Symbol} & \textbf{Meaning} \\
\midrule
$B$                 & the running BST (\figref{c05-running})   & $\mathit{lo},\mathit{hi}$ & open value bounds carried down a descent \\
$h$                 & height of the tree; every descent costs $\bigO{h}$ & $\mathrm{succ}(v)$ & next key larger than $v$ (inorder successor) \\
$\mathrm{in}(B)$    & inorder sequence of $B$ -- a sorted line  & $\mathrm{pre}(B)$ & preorder sequence (7.3's rebuild input) \\
$\mathcal{S}(v)$    & subtree rooted at $v$; its keys form an interval $[\min,\max]$ & $\lv{3}$ & difficulty badge (runs $\lv{1}$--$\lv{6}$) \\
\code{set}, \code{map} & \texttt{std} ordered containers -- self-balancing (Ch~8 proves them) & $\mathrm{lb}(x)$ & \code{lower\_bound}: first key $\ge x$ \\
\bottomrule
\end{tabular}
\end{center}
